\chapter{Introduction}\label{CH:introduction}
 Modern manufacturing relies heavily on producing precise, intricate metal components. This is often achieved through subtractive manufacturing techniques using equipment such as milling machines and lathes. These machines have been in use for approximately two hundred years \cite{groover_fundamentals_2010,moriwaki_multi-functional_2008}, and during this time have been continually improved. Manual control of these machines has been incrementally replaced by increasingly capable computerization \cite{reintjes_numerical_1991,groover_fundamentals_2010}. Modern milling machines and lathes often use Computer Numerical Control (CNC), which removes the need for a human machinist to control each detail of the machining process through manual inputs \cite{groover_fundamentals_2010}. Due to this continued computerization, these machines have grown in complexity, enabling more advanced machining operations, including advanced work with ductile materials, brittle materials, and increasingly complex part geometry \cite{moriwaki_multi-functional_2008,sheng_review_2026}. Continued development has also introduced new systems to these machines, including systems that allow the Machine to autonomously detect in real time a broken tool and replace the tool, detect abnormally high stress loads, or monitor motor loads across the machine. These innovations are significant because they enable the machines to autonomously catch in real time errors that previously would have been caught by a human, and with CNC machines could have propagated into damaged parts or damaged machines. These innovations work to reduce scrap part rates, damaged tooling, and critically machine downtime \cite{nagaraj_tool_2025,wong_technical_2020}.


\section{Machining Overview}\label{machiningOverview}
\subsection{Lathe}\label{Sect:lathe}
A lathe performs subtractive manufacturing by rotating the workpiece around the spindle axis, designated as the Z axis, while translating the cutting tool in the X and Z axes. The X axis is commonly used to move the tool radially and controls the finished diameter of the workpiece. The Z axis moves along the length of the workpiece and is used to control the axial extent of the cut. Given the workpiece rotates while the tool traverses a plane containing the axis of rotation, all features are produced as surfaces of revolution that are axially symmetric about the Z axis \cite{altintas_manufacturing_2012}. Lathes are commonly used to achieve part geometries that require turning outer diameters, boring internal diameters, facing, and drilling and tapping holes along the Z axis. Contoured profiles, including hemispherical and other curved forms are produced by coordinating the X and Z motions simultaneously \cite{groover_fundamentals_2010}. A conventional lathe has no Y axis; the Y axis is traditionally used only on mill-turn machines and multi-axis turning centers. The Y axis enables the off axis of rotation features which would break axial symmetry \cite{moriwaki_multi-functional_2008}.

During turning passes, the tool remains engaged with the workpiece continuously, and in the absence of material variation or progressive tool wear, the cutting force is approximately constant. This force is conventionally resolved into three orthogonal components, tangential, radial, and feed \cite{altintas_manufacturing_2012}. The magnitude of each of these forces is determined largely by the tool geometry, machine geometry, and the part characteristics. Large contributions to this are determined by the tool geometry and the lead and rake angles \cite{altintas_manufacturing_2012}. The three components are illustrated in Figure~\ref{fig:cutting_forces}.
%REVIEW-ADD-BEGIN
\begin{figure}[htbp]
	\centering
	\thesisfigure[0.9\linewidth]{ch1_cutting_force_components}%
		{Schematic of a turning tool engaged with a rotating workpiece, with arrows showing the tangential, radial, and feed components of the cutting force at the tool tip.}
	\caption[Cutting force components in turning]{Cutting force at the tool tip resolved into tangential, radial, and feed components.}
	\label{fig:cutting_forces}
\end{figure}
%REVIEW-ADD-END



\subsection{Mill}\label{Sect:mill}
A milling machine performs subtractive manufacturing by rotating the cutting tool about the spindle axis, designated the Z axis in a milling machine. Part geometry is achieved by moving the tool and workpiece relative to each other along three orthogonal directions. On a standard vertical machining center, the table carries the workpiece in the X and Y axes while the spindle head provides motion in the Z axis setting the depth of cut. Through these three degrees of freedom, material can be removed at arbitrary locations on the workpiece, so a far wider range of geometry can be produced compared to a lathe. Additional rotary axes extend this part geometry further. Four-axis machines add a single rotary axis about the X or Y axis and five-axis machines add a second rotary axis allowing the tool to approach the workpiece at a continuously variable orientation and make undercut and complex contoured surfaces without rotational symmetry possible  \cite{altintas_manufacturing_2012}.
The cutting force in a milling operation is applied to the workpiece through the flutes of the rotating tool, and each flute enters and leaves the cut once per spindle rotation. Engagement therefore is intermittent rather than continuous, and the load present on the machine is periodic at the tooth-passing frequency. The tooth-passing frequency is equal to the spindle rotational speed multiplied by the number of flutes on the cutting tool \cite{groover_fundamentals_2010,moriwaki_multi-functional_2008,altintas_manufacturing_2012}.

\subsection{CNC Based Manufacturing}\label{Sect:cnc}
Both lathes and mills can be improved by the addition of computer based control, and machines of this type are known as Computer Numerical Control (CNC) lathes and mills. Through CNC, tool paths are specified as coordinated motion commands rather than performed by an operator adjusting hand-wheels and cranks. CNC toolpaths allow for simultaneous coordinated multi-axis motions that are required for complex contoured geometry that would be impractical to produce manually \cite{groover_fundamentals_2010,altintas_manufacturing_2012}. On servo based machines, motion is regulated by nested closed-loop position and velocity controllers acting in each axis \cite{altintas_manufacturing_2012}. The controller compares the commanded position against measured position at each update of the servo loop, which often occurs faster than the machine could reach the desired location, position error is corrected continuously. This control system allows CNC machines to hold a high level of precision and to reproduce the same result across many parts. Machines driven by stepper motors are commonly used and typically operated in an open-loop configuration. This principle takes advantage of the internal stepping motion a stepper motor provides to achieve reproducible and precise control in a simplified method compared to servo based machines, trading some precision for lower cost \cite{henke_modeling_2013,ricci_simple_2018,binfet_implementation_2026}.

Adding sensing beyond the position control loop permits tighter control and the recording of process data during machining. This data can be retained on a per-part and per-machine basis and support process improvements across a production run in a manner that manual machining did not allow for \cite{tao_data-driven_2018,wong_technical_2020,altintas_manufacturing_2012, civiello_characterization_2026}.
High-performance CNC machines are typically expensive and occupy large floor footprints. These large often cast based footprints are used to raise the mass and structural damping of the machine, attenuating vibration and disturbance forces before tool disturbance can occur. This contributes substantially to the precision and repeatability these machines can achieve \cite{altintas_manufacturing_2012}. If these disturbances can instead be characterized and rejected within the control system, a portion of the precision obtained by this structural mass could in principle be obtained from control, allowing smaller and lighter machines to approach the performance of larger cast based machines. Machines of this type would be usable in applications where the machine must be enclosed and where cost or space required for this enclosure currently constrains what can be produced. 

\subsection{Observability of Machining Loads}\label{observability}
Some modern CNC machines are capable of recording substantial data during operation, and that data is presently used in three primary ways: to improve part quality across a production run, to reduce machine cycle time, and to indicate when maintenance should be performed \cite{tao_data-driven_2018,teti_advanced_2010}. Each of these acts after the fact, as the part being cut when the data was collected is already complete. With additional sensing and characterization of the machine, the same information could support correction during the current cut, acting on the workpiece currently being machined rather than the next one. Corrections of this kind would also accumulate, producing a record of which in-process adjustments succeeded which could be applied to a new part geometry without repeating the same trial cuts.

Systems performing monitoring of this kind are already commercially available. The Tool Monitoring Adaptive Control (TMAC), developed by Caron Engineering interfaces directly with a range of CNC controls, including Okuma platforms. The system monitors spindle motor power during the cut and operates on the principle that the power required to cut increases as the cutting edge wears. It applies this logic in two ways, it retires a worn tool or halts the machine on breakage, and in adaptive control mode, it overrides the programmed feed rate to hold spindle motor power near a set limit. Data from these cuts can be retained for later analysis and process improvements \cite{caron_tmac,cte_tmac_2020,okuma_caron, civiello_characterization_2026}.

The adaptive capability therefore exists, but it acts on an aggregate quantity. Spindle motor power represents the total power drawn by the cut. It does not resolve the force at the cutting edge into its principal components, and holding that power near a limit regulates the process only in the coarsest sense, without reference to the cutting conditions that have produced these changes. Civello \cite{civiello_characterization_2026} found that cutting force at the cutting edge could not be resolved from the data accurately enough to infer the resulting change in part geometry or surface finish. This limitation was attributed to the sample rate of the TMAC system, where the noise was no longer distinguishable from the signal. The signal was sampled at a rate of \SI{200}{\hertz}, which was too low to resolve the cutting-force signal above the noise. This resulted in the data that showed the process was feasible, however the electrical sampling system would need significant changes to be practical. A system with direct access to the motor’s electrical quantities and closer integration with the motor may resolve the signal-to-noise ratio problem and provide more insight than the TMAC system could. 


\subsection{Machine Load Transmission}\label{machineLoad}
When an axis is driven by a ballscrew, the cutting force acting on the carriage produces an axial load on the ballscrew. Given that a ballscrew is backdrivable the axial load appears at the motor as a rotational torque. The magnitude of this torque is set by the screw lead and the efficiency of the drive. The cutting load therefore reaches the motor through the existing drive system, without additional instrumentation added to the machine. 
%TODO Add figure that shows the load path from the cutting force
%REVIEW-ADD-BEGIN
\begin{figure}[htbp]
	\centering
	\thesisfigure[0.9\linewidth]{ch1_load_path}%
		{Block diagram tracing the cutting force on the workpiece through the machine carriage, the ballscrew nut, and the ballscrew to a torque on the motor shaft, and from that torque to the quadrature current drawn by the motor.}
	\caption[Load path from cutting force to motor current]{Load path from the cutting force at the tool to the quadrature current of the axis motor.}
	\label{fig:load_path}
\end{figure}
%REVIEW-ADD-END
 Many CNC machines come standard with ballscrews, however those that use alternative drive types such as lead screws can be retrofitted with ballscrews to achieve this behavior and enable this torque transmission. The resulting load path is shown in Figure~\ref{fig:load_path}.

Electric motors produce torque as a function of supplied current, the form of this current torque relationship varies on motor design. However, all electric motors fundamentally convert electrical current into a magnetic field that is used to produce torque. Some motor types produce torque at a fixed linear relationship to current, while others follow through a non-linear relationship dependent on rotor angle. Determining the load from the current therefore requires a priori knowledge of the electrical characteristics of the motor \cite{binfet_implementation_2026,derammelaere_linear_2010,henke_modeling_2013}. This torque to current relationship motivates the review of motor types in the following section.

Where the relationship is known, the cutting load appears at the axis drive as a disturbance torque, and the controller responds by changing the command current in order to preserve the axis trajectory. The cutting load is consequently present in the current signal and can be recovered from this disturbance. Recovery of this cutting load requires separating this load from the other torques the motor supplies during the same motion. These additional disturbance torques include friction in the system, inertial torque required to accelerate and decelerate the axis, and on a vertical axis the effects of gravity on the load of the carriage. Inertial and gravitational contributions follow from the axis mass and command trajectory can be predicted directly and therefore identified as a known disturbance. Friction is less tractable, comprising Coulomb, viscous, and velocity dependent components that vary with temperature and exhibit a hysteresis at velocity change \cite{altintas_manufacturing_2012,armstrong_survey_1994}. These disturbances are therefore estimated rather than predicted, and corrected for in situ through observer-based and signal processing techniques present in the proposed control system \cite{herbst_tuning_2023}.

The underlying principle for torque based disturbance characterization is already applied in a limited form on modern CNC machines. This is used where drive current detects events without any added sensor. These methods are commonly described as “sensorless”. The term originates in motor control, where it denotes the estimation of rotor position without a position sensor \cite{cho_tool_2016,derammelaere_linear_2010}. In this context it is used in the broader sense adopted in machine tool practice, meaning that the quantity of interest is inferred from signals already present in the drive rather than measured by a sensor added for that purpose. Homing against a hard stop is one example. When the axis reaches the stop, the motion ceases while the controller continues to the command position. This position error accumulates and the commanded current rises toward the motor limit. This sustained rise identifies a hard stop and can be used to establish a machine reference point without a limit switch \cite{adi_tmc2209_2023,voron_sensorless}. This hard stop referencing method is only practical due to the relatively high speeds at which the controller is able to sample error and detect a disturbance so no damage to the machine occurs during this brief period of time when the hard stop is in contact with the carriage on that axis.

This principle can be expanded to detect other machine changes, tool breakage is another common use case for detecting sudden current load changes. This is normally performed by monitoring the spindle drive rather than an axis drive. The sudden loss of a cutting edge produces an abrupt decrease in spindle load, most commonly a decrease in power as the tool is no longer in contact with the workpiece \cite{cho_tool_2016,wong_technical_2020}. This signal is drawn from a different motor, however the principle is the same in both use cases. Both the axial hard stop and the tool breakage are determined by sudden change in current draw in the motor controlling that motion. The present work extends this same principle from an event detection approach to a continuous estimation of the load magnitude. 

\subsection{Characterizing Load at the Cutting Edge}\label{characterizingLoad}
With a motor and control system of sufficient electrical and control bandwidth, current can be sampled rapidly enough to follow variation in the cutting load rather than only tracking an average value. Turning is the appropriate operation in which to establish this relationship first. The tool remains engaged continuously through a turning pass, and during a hemispherical operation translates in both the X and Z axes causing a change in the force vector direction. With these passes, in the absence of material variation or progressive tool wear, the cutting force is approximately steady and the motor torque is likewise steady. Milling by contrast, presents an intermittent load at the tooth-passing frequency \cite{altintas_manufacturing_2012}, superimposing a large periodic component on the signal of interest. Beginning with turning isolates the current-to-force relationship from the periodic relationship which allows it to be established under the simplest available loading condition. The two loading conditions are compared in Figure~\ref{fig:turning_vs_milling}.

%TODO Add figure of turning force versus milling force - show tooth-passing frequency if possible. 
%REVIEW-ADD-BEGIN
\begin{figure}[htbp]
	\centering
	\thesisfigure[0.9\linewidth]{ch1_turning_vs_milling_force}%
		{Two plots of cutting force versus time: a nearly constant force for a turning pass, and a force for a milling pass that rises and falls once per tooth, repeating at the tooth-passing frequency.}
	\caption[Cutting force in turning and milling]{Cutting force during turning, which is approximately constant, and during milling, which is periodic at the tooth-passing frequency.}
	\label{fig:turning_vs_milling}
\end{figure}
%REVIEW-ADD-END

Direct instrumentation of the cutting zone is not viable in the target applications, where the machining environment is enclosed and access is limited. Introducing sensors into this environment introduces a non-trivial cost aspect and would expose sensors to potential damage due to conditions internal to the machine during operation\cite{wong_technical_2020}. Using already existing mechanics to infer the load from the motor current instead, locates the measurement at the drive, where the necessary electrical quantities are already present in the controller, and requires no instrumentation in the cutting zone at all. Establishing this relationship also opens a path to a quantity that is not otherwise observable during machining. Force at the cutting edge is distributed over the contact between the tool and the workpiece, and the pressure in that contact, rather than the total force, is what governs the material’s response to the cut. Pressure follows from force and the area over which it acts, and that area is determined by the chip cross-section and geometry of the cutting tool \cite{altintas_manufacturing_2012, civiello_characterization_2026}.

 Specific cutting energy provides the standard basis for relating the measured force to the condition in the cutting zone \cite{altintas_manufacturing_2012, groover_fundamentals_2010, civiello_characterization_2026}. Where force can be recovered from drive current, tool geometry, and chip cross-section from machining parameters, an estimate of pressure at the cutting edge can be determined without additional measurement.
 
A pressure estimate obtained through this could be of use where a material’s machining envelope depends on the loading conditions it experiences rather than on nominal properties alone. Certain material classes undergo solid-state phase transforms within the range of stress and temperature encountered at the cutting edge, and for these deformations the behavior is dependent on the pressure and temperature history the material has experienced \cite{woodrum_developing_2017}. Machining such materials predictably requires a priori knowledge of the load actually being applied at the cutting edge, which standard commercial CNC control systems do not provide \cite{teti_advanced_2010,woodrum_developing_2017}.


Recovering pressure in this manner is beyond the scope of the present work, which establishes a relationship between motor current and axis load that any such estimate would depend upon. This approach is identified here because it defines the direction in which this capability could extend, and because the requirements this imposes depend on the method of data collection in the present work. These requirements are bandwidth sufficient to follow the cut, a torque-to-current relationship, and a control system capable of accepting current feedback. Meeting these requirements on a stepper-driven axis requires a motor model, a control architecture capable of maintaining the torque-to-current relationship under load, and a processing platform capable of operating at the required rate. Each of these is developed in the present work. 

\section{Electrical Overview}\label{sect:electricalOverview}
To accurately capture the present work, the fundamentals of common motor types are reviewed. Motor types determine the torque-to-current relationship. Direct Current (DC) motors have a relationship such that the torque produced is directly proportional to the current passing through the motor windings. The torque to current relationship in a DC motor does not depend on the motor rotor angle for torque generation and is constant regardless of the current rotor angle. Brushless Direct Current (BLDC) motors and Hybrid Stepper Motors do depend on the rotor angle as part of the current to torque output relationship along with the current moving through the motor windings \cite{binfet_implementation_2026,nilsson_universal_2013}.

\subsection{DC Motor}\label{dcMotors}
DC Motors have the most direct relationship for current and torque. In a DC motor, a voltage is applied to the motor windings, which are located on the rotor, and permanent magnets inside on the stator respond to the generated magnetic field and rotation is generated. While this motor type provides a simple installation and simple operation, the mechanical link of using brushes to pass the current from the fixed motor connections through the commutator to the rotating rotor is susceptible to failure from wear and can lead to a shortened motor life compared to technologies such as BLDC motors \cite{binfet_implementation_2026,nilsson_universal_2013}. The construction is shown in Figure~\ref{fig:dc_motor}.

%TODO (Insert figure of the internals of a regular DC Motor)
%REVIEW-ADD-BEGIN
\begin{figure}[htbp]
	\centering
	\thesisfigure[0.6\linewidth]{ch1_dc_motor}%
		{Cutaway of a brushed DC motor showing the permanent magnets on the stator, the wound rotor, and the brushes contacting the commutator.}
	\caption[Brushed DC motor construction]{Internal construction of a brushed DC motor.}
	\label{fig:dc_motor}
\end{figure}
%REVIEW-ADD-END

\subsection{Brushless DC Motors}\label{brushlessMotors}
Brushless DC Motors are often used in servo-based CNC machines where the longevity and precision control of the axes are required. In these applications a BLDC motor is selected due to the lack of brushes internal to the motor that can degrade with use. With no direct contact between the rotor and the stator, BLDC motors are not susceptible to the same wear and degradation as a DC motor with brushes in the commutator failing. Therefore BLDC motors have less wear internally that cause premature motor failure. These motors require additional control to operate as these motors require timing to three separate internal phases to properly generate rotational motion. BLDC motors differ from Permanent Magnet Synchronous Motors (PMSM) in the way that the motors are driven and the torque output relationship. In PMSM motors the coil windings are such that an AC signal is required to drive the motor, meaning external sinusoidal control of the voltage is needed. Whereas with the BLDC motor, the motion can be generated with a trapezoidal method instead of a sinusoidal method \cite{binfet_implementation_2026,nilsson_universal_2013}. The construction is shown in Figure~\ref{fig:bldc_motor}. 

%TODO (Insert figure of the internals of a brushless DC motor)
%REVIEW-ADD-BEGIN
\begin{figure}[htbp]
	\centering
	\thesisfigure[0.6\linewidth]{ch1_bldc_motor}%
		{Cross-section of a three-phase brushless DC motor showing the stator windings for phases A, B, and C and the permanent-magnet rotor.}
	\caption[Brushless DC motor construction]{Internal construction of a three-phase brushless DC motor.}
	\label{fig:bldc_motor}
\end{figure}
%REVIEW-ADD-END

\subsection{Hybrid Stepper Motors}\label{hybridStepperMotor}
Hybrid Stepper Motors are used when repeatable motion is desired at a low cost. The hybrid stepper motor takes advantage of internal fixed coils that cause predictable and repeatable rotation in a step motion. The internal rotor can be thought of as moving akin to a gear where each coil will advance the gear by one quarter of a tooth. While this is not the direct method used in hybrid stepper motors, this model provides a more clear understanding of the intended use case for the stepper motor. Internal to a stepper motor the rotor and stator poles are set such that both magnetically align at each step at the maximum magnetic field strength. This electrical and magnetic coupling is what provides the strong holding force of a stepper motor along with the stepping action \cite{henke_modeling_2013,binfet_implementation_2026,derammelaere_linear_2010}. A standard method to driving a stepper motor is to use pulse trains alternating between the two phases of the motor. This electrical signal is simple to produce and causes rotational motion of a stepper motor in a way that can be easily discretized \cite{ricci_simple_2018,lai_modelling_2018}. The construction is shown in Figure~\ref{fig:hybrid_stepper}. 

%TODO (Insert figure of Hybrid Stepper Motor)
%REVIEW-ADD-BEGIN
\begin{figure}[htbp]
	\centering
	\thesisfigure[0.6\linewidth]{ch1_hybrid_stepper_motor}%
		{Cross-section of a two-phase hybrid stepper motor showing the toothed stator poles carrying the phase A and phase B windings and the toothed permanent-magnet rotor.}
	\caption[Hybrid stepper motor construction]{Internal construction of a two-phase hybrid stepper motor.}
	\label{fig:hybrid_stepper}
\end{figure}
%REVIEW-ADD-END

\subsection{Field Oriented Control}\label{FOC}
Field Oriented Control (FOC) is commonly used on BLDC motors and three-phase power systems to transform the multiple phases into two orthogonal components, the quadrature ($q$) axis and the direct ($d$) axis. The quadrature axis is the torque-generating axis that is always leading the electrical angle by 90 degrees while the direct axis is aligned with the rotor flux and therefore is not producing a rotational torque. This method is commonly used on BLDC motors as FOC reduces the complexity of driving the BLDC motor down to two vectors, the Quadrature which is scaled by the required torque output and the Direct which is minimized. Through this transform, the motors are able to maximize the torque produced and linearize the torque-to-current relationship \cite{benchaib_discrete-time_2003,nilsson_universal_2013,binfet_implementation_2026}. The transforms are illustrated in Figure~\ref{fig:park_bldc}.

%TODO (Include figure of BLDC through Park transform)
%REVIEW-ADD-BEGIN
\begin{figure}[htbp]
	\centering
	\thesisfigure[0.6\linewidth]{ch1_clarke_park_three_phase}%
		{Vector diagram showing the three stationary phase axes a, b, and c at 120 degrees, the stationary alpha-beta frame produced by the Clarke transform, and the d-q frame rotating at the electrical angle theta e produced by the Park transform.}
	\caption[Clarke and Park transforms for a three-phase motor]{Clarke and Park transforms for a three-phase motor, from the stationary $a$--$b$--$c$ phase axes to the rotating $d$--$q$ frame.}
	\label{fig:park_bldc}
\end{figure}
%REVIEW-ADD-END

The present work seeks to show that this relationship can be extended from the BLDC motor to the Hybrid Stepper Motor \cite{lai_modelling_2018,ricci_simple_2018,bernardi_performance_2022,derammelaere_linear_2010}. The same fundamental property is being used to convert the discrete steps internal to the stepper motor into a continuously rotating system. The Quadrature current is the current solely responsible for generating rotational torque in the stepper motor and the direct-axis current is wasted energy that does not contribute to the rotation or stepping motion of the motor. This is because the direct-axis current generates a flux along the rotor magnet axis which does not generate torque but still passes current through the motor. Through this transform, the Hybrid Stepper Motor will be treated as a synchronous machine where two continuous sine waves can be applied to the two phases of the stepper motor at 90 degrees out of phase to generate rotation. The phase of both sine waves is directly linked to the motor rotor electrical angle keeping both in phase with the motor to generate the rotation \cite{derammelaere_linear_2010,bernardi_performance_2022,lai_modelling_2018}. Figure~\ref{fig:two_phase_dq} shows the two-phase case. 

%TODO (Include figure of Clarke Transform and Stepper Motor) A two-phase machine does not need a Clarke transform. Label this as two-phase $a$--$b$ / $d$--$q$ diagram.
%REVIEW-ADD-BEGIN
\begin{figure}[htbp]
	\centering
	\thesisfigure[0.6\linewidth]{ch1_park_two_phase}%
		{Vector diagram of a two-phase hybrid stepper motor showing the orthogonal phase A and phase B axes, and the d and q axes rotating at the electrical angle theta e, with the current vector aligned to the q axis.}
	\caption[Park transform for a two-phase stepper motor]{Park transform for a two-phase hybrid stepper motor. The phase axes are already orthogonal, so no Clarke transform is needed.}
	\label{fig:two_phase_dq}
\end{figure}
%REVIEW-ADD-END

\section{Control Systems Overview}\label{controlSystemOverview}


The current work covers two control systems which necessitates a review of the basics of both control systems. Traditional control systems are computed and designed in the frequency domain while modern control systems are often completed in the time domain through the use of linear algebra \cite{astrom_feedback_2021,binfet_implementation_2026}. The initial design for implementing FOC control on a stepper motor was achieved through a cascaded PID control loop. This provided some initial results and proved the system was worth pursuing further. One of the limitations of this solution came from the hardware running it, a Raspberry Pi 4. This introduced inconsistencies in timing and created a non-deterministic control system. Due to these computational limits, the system was moved to an FPGA which also provided the possibility of doing computation using matrices while having deterministic timing. With this, the present work uses an Active Disturbance Rejection Control method to remove additional disturbances in the system through the use of observers. Despite the advantage of having matrices to represent the control law of the system, a minimum footprint Active Disturbance Rejection Control algorithm was used. This converted the matrix mathematics to a transfer function that requires fewer FPGA resources \cite{binfet_implementation_2026,herbst_minimum_2021}. 

\subsection{Proportional Integral Derivative Control}\label{pidOverview}
One classical approach to control systems for CNC machines is to use Proportional Integral Derivative (PID) control to control the motion of the CNC machine. A PID system makes use of comparing the feedback to the target reference point to calculate the error. This traditional control system is often implemented as a single input single output system (SISO) where the position and velocity control are handled separately\cite{altintas_manufacturing_2012,binfet_implementation_2026}.


%REVIEW-ADD-BEGIN
The PID control law is
\begin{equation}
	u(t) = K_p\, e(t) + K_i \int_{0}^{t} e(\tau)\, d\tau + K_d\, \frac{de(t)}{dt},
	\qquad e(t) = r(t) - y(t),
	\label{eq:pid}
\end{equation}
where $r$ is the reference, $y$ is the measured output, and $K_p$, $K_i$, and $K_d$ are the proportional, integral, and derivative gains \cite{astrom_feedback_2021}. The structure of the controller is shown in Figure~\ref{fig:pid_block}.
%REVIEW-ADD-END
%REVIEW-ADD-BEGIN
\begin{figure}[htbp]
	\centering
	\thesisfigure[0.9\linewidth]{ch1_pid_block}%
		{Block diagram of a PID controller: the error between the reference and the measured output feeds parallel proportional, integral, and derivative paths whose sum is the control signal applied to the plant.}
	\caption[PID controller]{Parallel PID controller acting on the error between the reference and the measured output, Equation~\eqref{eq:pid}.}
	\label{fig:pid_block}
\end{figure}
%REVIEW-ADD-END

Through this a position motion profile can be created where the velocity control system is quickly saturated to the reference feedrate specified in the G-Code for the program. Through this, a responsive system that is able to reach the target position which a high level of accuracy is achieved while maintaining the desired feedrate on the machine. This modification of a traditional position control system allows for a set velocity target to be maintained during cutting operations which is vital to the proper functionality of a milling machine or lathe. Constant cutting velocity provides a consistent chip load which produces a more precise final result \cite{altintas_manufacturing_2012}.


%REVIEW-ADD-BEGIN
\begin{figure}[htbp]
	\centering
	\thesisfigure[0.9\linewidth]{ch1_cascaded_position_velocity}%
		{Block diagram of cascaded position and velocity loops: the position controller output is the velocity reference, the velocity controller output is the current or torque command, and position and velocity feedback return from the motor encoder.}
	\caption[Cascaded position and velocity control]{Cascaded position and velocity control, in which the position controller supplies the velocity reference.}
	\label{fig:cascaded_pid}
\end{figure}
%REVIEW-ADD-END

To achieve FOC control using PID systems, a preliminary design was created that utilized a PID controller for the position, velocity, quadrature current, and direct current. This system involved cascading the control system such that the output of one controller became the reference for the next controller. Through this, a reference quadrature current would be established that would provide the correct current to generate the target velocity and move the motor to the desired location at the desired velocity. By using the Park Transform, the quadrature current was made linear in respect to the output torque \cite{benchaib_discrete-time_2003,derammelaere_linear_2010}. The quadrature current reference is treated as a step target with a constant current being related to a constant acceleration while under no-load operation. During a constant load, instead of the quadrature current settling at zero with no motion, the quadrature current will settle at a value directly proportional to the load on the system. The current loops are shown in Figure~\ref{fig:pid_current_loop}. Running in parallel to the quadrature current control loop was the direct current control loop with a constant reference input of zero. Through this, a linear current target change can be attributed to disturbances in the motion of the motor such as an increase of torque working against the rotation of the stepper motor.


%REVIEW-ADD-BEGIN
\begin{figure}[htbp]
	\centering
	\thesisfigure[0.9\linewidth]{ch1_pid_current_loops}%
		{Block diagram of the preliminary PID current control: measured phase currents pass through the Park transform to give i d and i q, two PID controllers regulate i d to zero and i q to its reference, and the inverse Park transform produces the phase voltage commands.}
	\caption[Preliminary PID current control]{Preliminary PID current control in the rotor-fixed $d$--$q$ frame.}
	\label{fig:pid_current_loop}
\end{figure}
%REVIEW-ADD-END

Early success was found with this system, however this was implemented on a Raspberry Pi 4 B 8 Gb model which introduced limitations in how fast the control system could update and how fast feedback could be collected. This limited the control system to updating at approximately \SI{200}{\hertz}. At \SI{200}{\hertz} the update period is \SI{5}{\milli\second}, which is longer than the \SI{3}{\milli\second} electrical time constant of the motor windings. This limitation meant the time constant of the stepper motor was shorter than the update period of the control system. Through this, no meaningful change in current could be detected that could be related to the change in torque observed. However, this initial system did prove that a stepper motor could be run using FOC and that a change in quadrature current was noticeable in the offset of the waveform as the load increased. Given limitations in the software system on the Raspberry Pi and potential additional disturbances in the mechanical system imperfections, a conclusive relationship between torque and current could not be determined. However, the initial data showed an increased trend in quadrature current when torque increased. 

\subsection{Active Disturbance Rejection Control}\label{ADRCOverview}
A principal limitation of PID control is the delay in detecting disturbances through the error signal from the plant. With PID control, there is no predictive work being done to estimate the error that will result from the plant and how to separate out error from the plant versus error that is the result of a disturbance. With the Active Disturbance Rejection Control (ADRC) the error of the plant is predicted through an extended state observer (ESO). Modeling this error using the current reading from the plant allows for converging the estimate of the plant response to the control signal. Through this, the way the plant will react to a given control signal can be efficiently modeled and disturbances to this control signal predicted \cite{herbst_tuning_2023}. 
By using an ADRC system, the disturbances of the plant can be modeled and removed which allows for a clean control signal to be provided to the plant. This results in a quadrature current reference that is more predictable and will account for external loads that were not modeled in the stepper motor \cite{herbst_tuning_2023}. To achieve this result, a minimum footprint ADRC system was modeled in Simulink where an ADRC is used to control the velocity of the motor and the current in both the quadrature and direct phase. A proportional controller is used to control the position of the motor, as the output of the position controller following convention should saturate to the velocity limit imposed by the machining operation. The structure of the ADRC is shown in Figure~\ref{fig:adrc_block}. 

%TODO FIgure ADRC block diagram showing the plant, the extended state observer, and the control law (Chapter 2, Eqs.~(eso) and (adrc\_law)).}
%REVIEW-ADD-BEGIN
\begin{figure}[htbp]
	\centering
	\thesisfigure[0.9\linewidth]{ch1_adrc_block}%
		{Block diagram of first-order ADRC: the plant output y feeds an extended state observer that estimates y and the total disturbance f; the control law subtracts the disturbance estimate from a proportional term and divides by b0 to form the control signal u applied to the plant.}
	\caption[First-order ADRC structure]{First-order ADRC. The extended state observer estimates the output and the total disturbance, and the control law cancels the estimated disturbance (Equations~\eqref{eq:eso} and~\eqref{eq:adrc_law}).}
	\label{fig:adrc_block}
\end{figure}
%REVIEW-ADD-END


\section{Hardware Platform Overview}\label{FPGAOverview}

A Field Programmable Gate Array (FPGA) is an integrated circuit made up of an array of configurable logic blocks which are connected through a programmable routing network. Each logic block contains lookup tables (LUTs) that can implement combinational logic functions, along with flip-flops that hold the state of the circuit between clock cycles. Many FPGAs also include dedicated resources such as Digital Signal Processing (DSP) slices for fast multiplication and accumulation, block RAM for on-chip memory storage, and configurable input and output blocks for interfacing with external hardware. The behavior of the FPGA is described using a Hardware Description Language (HDL), such as VHDL or Verilog. This description is synthesized into a netlist of logic elements, placed and routed onto the physical resources of the device, and then loaded onto the FPGA as a configuration bitstream \cite{binfet_implementation_2026}. A simplified view of this architecture is shown in Figure~\ref{fig:fpga_architecture}. The key difference between an FPGA and a traditional processor such as the Raspberry Pi 4 is in how an algorithm is executed. A microprocessor executes a sequence of instructions in a serial fashion with dedicated instructions to use, while an FPGA implements the algorithm directly as digital circuitry. Through this, independent operations such as the Clarke and Park transforms, the extended state observers, and the encoder decoding can all be computed at the same time in separate regions of the device rather than competing for time on a single processor.

%TODO (Insert figure of the internal architecture of an FPGA showing logic blocks, routing, and I/O blocks)
%REVIEW-ADD-BEGIN
\begin{figure}[htbp]
	\centering
	\thesisfigure[0.7\linewidth]{ch1_fpga_architecture}%
		{Diagram of FPGA internal architecture: a grid of configurable logic blocks containing lookup tables and flip-flops, connected by programmable routing, with columns of DSP slices and block RAM and input/output blocks around the edge.}
	\caption[FPGA architecture]{Simplified internal architecture of an FPGA.}
	\label{fig:fpga_architecture}
\end{figure}
%REVIEW-ADD-END

A System on Module (SoM) FPGA packages the FPGA together with the supporting hardware required for operation onto a single compact circuit board. This supporting hardware typically includes the power regulation and power sequencing, the clock sources, external DDR memory, and non-volatile flash memory for storing the configuration bitstream. The SoM is mounted onto a carrier board through high-density connectors, where the carrier board provides the application-specific hardware such as the motor drivers, the current sensing circuitry, and the encoder connections. Many SoMs are built around an SoC FPGA, which combines hard processor cores and programmable logic on the same chip. In these devices, the processing system can run an operating system to handle tasks that are not time critical, such as parsing G-Code, communicating with a user interface, and logging data, while the programmable logic handles the time critical control loops. The two sides of the device communicate through an on-chip bus such as the Advanced eXtensible Interface (AXI) \cite{binfet_implementation_2026}. The advantage of using a SoM is that the most difficult portions of the board design, such as the high-speed memory routing and the power sequencing of the FPGA, have already been designed and validated by the manufacturer. Therefore, the design effort can be focused on the motor control hardware and the control system itself rather than on the supporting electronics of the FPGA.


The control system in the present work was designed in Simulink with the intent of being converted to HDL for synthesis. While Simulink allows a control system to be modeled with very few restrictions, a model intended for HDL synthesis must describe something that can be feasibly built out of physical digital logic, which imposes several constraints on the design. The model must be in discrete time with a fixed-step solver, so continuous blocks such as integrators and transfer functions are replaced with discrete equivalents. The transfer function form of the minimum footprint ADRC is well suited to this, as a discretized transfer function maps directly to a set of delays, multipliers, and adders. The model is also converted to fixed-point arithmetic, as floating-point arithmetic consumes a large amount of FPGA resources, which requires a word length and fraction length to be selected for each signal to balance range, resolution, and resource usage \cite{binfet_implementation_2026}. Only blocks with HDL support may be used, so functions such as the sine and cosine required for the Park transform are implemented using lookup tables. Additionally, long chains of arithmetic operations must be pipelined to meet the timing of the FPGA clock, and the resulting latency must be included in the simulation so that the behavior of the model matches the behavior of the hardware. The resulting partition of the model is shown in Figure~\ref{fig:simulink_partition}.

%TODO (Insert figure of the Simulink model showing the partition between the HDL subsystem and the plant model)
%REVIEW-ADD-BEGIN
\begin{figure}[htbp]
	\centering
	\thesisfigure[0.9\linewidth]{ch1_simulink_partition}%
		{Top-level diagram of the Simulink model divided into the HDL-synthesizable controller subsystem and the simulation-only subsystem containing the stepper motor and PCB models, with the PWM, ADC, and encoder signals crossing the boundary.}
	\caption[Partition of the Simulink model]{Partition of the Simulink model into the HDL-synthesizable controller and the simulation-only plant.}
	\label{fig:simulink_partition}
\end{figure}
%REVIEW-ADD-END

The largest benefit of moving the control system to an FPGA is deterministic timing. On a general purpose processor running an operating system, as with the Raspberry Pi 4 used in the initial design, the time between control loop updates varies due to the operating system scheduler, interrupts, other running processes, and memory access delays. A control system designed in discrete time assumes that the sample period is constant, and the integral terms, derivative terms, and extended state observer are all computed using this sample period. When the actual period differs from the assumed period, the integration accumulates error, the derivative estimate becomes noisy, and the observer estimate deviates from the true state of the plant. This timing jitter also acts as a varying delay in the feedback path which reduces the phase margin of the control system \cite{binfet_implementation_2026, cervin_how_2003, franklin_digital_1998}.

On an FPGA, each control loop update completes in a fixed number of clock cycles derived from the system clock. Through this, the sample period is constant to within the accuracy of the clock, and the latency from a measurement to the corresponding control output is known exactly and can be included in the design of the control system. Because each operation runs in its own dedicated logic, adding computation or an additional motor axis does not affect the timing of the existing control loops. The sampling of the phase currents can also be synchronized with the PWM signals driving the motor, such that the current is measured at the same point in every PWM cycle. This isolates the sampling from the switching noise of the motor driver \cite{blasko_sampling_1998}. Together, these properties allow the control loop to be updated at a rate many times faster than the approximately \SI{200}{\hertz} achieved on the Raspberry Pi 4, and faster than the electrical time constant of the stepper motor windings. In the present design the current loops update at \SI{50}{\kilo\hertz}, 250 times the \SI{200}{\hertz} rate of the Raspberry Pi 4 system. The update rates of both implementations are compared in Table~\ref{tab:update_rates}

%\note{If you measured the loop-time jitter on the Raspberry Pi, add it to the Timing column.}
%TODO (Insert achieved control loop update rate and a comparison of the timing jitter between the Raspberry Pi 4 and the FPGA)
%REVIEW-ADD-BEGIN
\begin{table}[htbp]
	\centering
	\caption{Control loop update rates of the Raspberry Pi 4 and FPGA implementations.}
	\label{tab:update_rates}
	\begin{tabular}{@{}ll>{\raggedright\arraybackslash}p{0.17\linewidth}>{\raggedright\arraybackslash}p{0.3\linewidth}@{}}
		\toprule
		Implementation & Function & Update rate & Timing \\
		\midrule
		Raspberry Pi 4 & All loops (cascaded PID) & $\approx\SI{200}{\hertz}$ & Varies with operating-system scheduling \\
		FPGA & Current loops & \SI{50}{\kilo\hertz} & Fixed, 4000 clock cycles \\
		FPGA & Velocity loop & \SI{2.5}{\kilo\hertz} & Fixed, 80000 clock cycles \\
		FPGA & Feedback processing, PWM & \SI{200}{\mega\hertz} & Every clock cycle \\
		\bottomrule
	\end{tabular}
\end{table}
%REVIEW-ADD-END